\section{Source Section III.E: Steady Two-Phase Heat Conduction}
\label{sec:steady-heat-conduction}

\sourceeqrange{Eqs. (3.17)--(3.19) and figures 8--10}

\subsection*{Purpose}

This subsection gives a compact definition and units guide for the steady
heat-flow scenario.  It does not verify the reported Nusselt values.

\subsection*{Reading Guide}

The definitions below identify which scalar outputs the source reports and
which finite checks the companion can safely perform.

\subsection*{Effective Conductivity and Heat Flux}

Equation (3.17) defines an effective thermal conductivity from the bottom-wall
heat input and the imposed temperature difference.  Equation (3.18) defines the
local bottom heat flux from the conductive heat flux at the wall where the
velocity vanishes.  Equation (3.19) forms a dimensionless heat-transport ratio by
normalizing the effective conductivity by a liquid reference conductivity.

The finite bookkeeping chain is:
\begin{align}
  \lambda_{\rm eff}
  &=
  \frac{\displaystyle\int_0^L Q_b(x)\,\dd x}{T_b-T_t},
  \sourceeqbadge{3.17}
  \label{eq:iii-effective-conductivity-ledger}\\
  Q_b(x)
  &=-\left.\lambda\,\partial_yT\right|_{y=0},
  \sourceeqbadge{3.18}
  \label{eq:iii-bottom-heat-flux-ledger}\\
  Nu&=\frac{\lambda_{\rm eff}}{\lambda_\ell}.
  \sourceeqbadge{3.19}
  \label{eq:iii-nusselt-ledger}
\end{align}
Equation~\eqref{eq:iii-effective-conductivity-ledger} is an integrated
definition over the bottom boundary, not a local differential law.  The
conductive expression in \eqref{eq:iii-bottom-heat-flux-ledger} is the
source-stated Fourier-type wall heat-flux diagnostic, evaluated on the wall
branch where the velocity vanishes.  The ratio
\eqref{eq:iii-nusselt-ledger} is dimensionless by construction; reported values
of \(Nu\) remain source simulation outputs.

The useful bookkeeping point is the order of operations: compute the wall
heat-flux diagnostic on the boundary branch, integrate it along the bottom wall,
divide
by the imposed temperature difference, and only then normalize by the liquid
reference conductivity.  Reversing this order would confuse a local flux with
the integrated effective-conductivity definition.


The source also gives a one-phase no-convection estimate after Eq. (3.19).  In
that branch the effective conductivity reduces to the local conductive
coefficient.  Using Eq. (3.4),
\begin{equation}
  \lambda=k_B\nu_0n,
  \qquad
  \lambda_\ell=k_B\nu_0(1.7n_c),
\end{equation}
so
\begin{equation}
  Nu=\frac{\lambda}{\lambda_\ell}
  =\frac{k_B\nu_0n}{k_B\nu_0(1.7n_c)}
  =\frac{n}{1.7n_c}.
  \sourceeqbadge{3.19}
  \label{eq:iii-one-phase-nusselt-estimate}
\end{equation}
This estimate is a reference branch for reading the plotted Nusselt values; it
is not a calculation of the two-phase convective state.

\subsection*{Finite Checks}

The verification layer checks that the effective-conductivity definition has the
same dimensions as a thermal conductivity in the source's two-dimensional cell
bookkeeping and that \(Nu=\lambda_{\rm eff}/\lambda_\ell\) is dimensionless.  It
also records that any measured value of \(Nu\) remains a simulation output.

The finite check is therefore dimensional: a ratio of like conductivities is
dimensionless.  No steady heat-conduction solution or reported transport value
is recomputed here.
The effective-conductivity, bottom-heat-flux, and Nusselt-ratio checks are
covered by the Section III mirror pair listed in
Section~\ref{sec:numerical-results-overview}.  These checks do not solve the
steady two-phase heat-conduction problem.

\subsection*{Limitations}

The presence or absence of gravity changes the state being read, but the guide
only records the branch conditions.  It does not calculate the steady state or
validate the heat-pipe-like transport mechanism.
